\daymark{29}{the Last Day, Y.H. 812}

The bread went into the oven before the light had cleared the sill, and the house took the heat the way a vessel takes water that is poured too fast --- carefully, and with more of it going than staying. \person{Corin} was carried close enough to feel it, in the basket at \person{Sella Vail}'s hip under \place{the Vail kiln-house} roof, and the warmth came to him in levels: first the oven mouth when she drew the lid, then the standing heat in the room, then the thin edge of it that reached the fold at his chin and went no further. He had spent every day of his life in rooms where heat was held in something and counted as it went. A man who had once put both hands inside a stranger's chest to keep a failing heat from spilling had no hand now for the loaves going in beside him. The bread was the last bread of the year. They baked it as they baked every loaf, without a word about the year.

The loaves came out set a little apart from one another, brown-edged and set firm, and Sella moved her lips over them as she set them on the rack. It was a counting that kept no figure. Three here, two there, one turned with its brown edge up and looked at twice, the same arithmetic she had been moving along the shelf for a tend and more, so that no number need ever be said aloud where a slate might be near. The smell of barley smoke went into the wet wool and stayed there. The boy lay with the fold at his ear and let the heat come and go off the loaves in small waves, each wave weaker than the one before, and thought of the way a room cools after the fire is banked and before it settles. He knew the order of that cooling the way other men know the order of their own stairs. It did not help him. There was still no hand in the house that would take him the first pace toward the rack.

In the fold the two crusts lay where they had been put. The older was stiff past the smell of anything, brown-edged and cold since Stone 15; the newer still gave a thin breath of barley smoke when the wind turned off the brook, though it had been given on Stone 28 and had no warmth left in it for a week. Neither had been lifted. Neither had been named. No third crust came into the fold on \lore{the Last Day}, though bread was going out of the oven in armfuls and the house was full of the smell of it. He understood the rule by which that happened without being told: a gift that cannot be answered cannot be refused, and a gift that cannot be refused is best not spoken of, because a thing spoken of becomes a thing a person might have to answer. So the loaves went on the rack and the crusts stayed in the fold, and the two records of the same barley lay a hand's breadth apart and were never put together.

Bread. Two crusts. No third.

\scenebreak

\person{Mira Silt} came up from the eels in the middle of the morning with pale slime dried on her sleeves and silt in the creases of her neck, her head held a little to the right against the cold in the place where the top of her left ear had been. She had come up the yard from \place{the Fishstairs} and stopped where the warmth came out of the door and went no farther, as though she had meant to pass it and thought better at the last pace. A woman at the tubs asked her whose bread was going past, meaning the smell of it, or meaning the basket at the hip; Mira answered the question with a question and her own answer before the woman could choose between them.

“Whose do you think?” she said. “It came from a warm house, that's all. It's bread, and it'll be eaten, and that's the whole of it. You want to know whose, and you don't want to know. I can see which one's winning.”

She flicked an eel-bone against her thumb while she said it, a small grey thing from the pocket of her apron, and put it away when the deciding was done. The small smile stayed on her face the way a door stays on the latch, which can be shut without anybody noticing the movement. She did not look into the basket. She had not once looked into it all the time she had stood in the warmth.

Someone at the tubs said a figure then --- how many loaves had gone into the house for the feast, or how many the moot would take, or what a pair of eels had fetched before \lore{the Five} --- and Mira's thumb went under her apron to the grey wool cord at her waist and stayed there while she decided what she had heard. The boy watched the hand go and watched it come back out with another bone between finger and thumb. When she shifted her weight to go back down the yard, a second cord twisted twice beneath the grey showed for a moment below her belt and was gone under the apron again: the one that had never been carried up the river to \place{Nine Pots} and never lifted at a door. The knots on neither were counted. She kept nothing on slate, this woman, and what she kept under her apron had no name written on it and no figure anybody could take from her. He put the hand on one place of the tray and the two cords on another and did not make a third thing out of their being the same colour, because a likeness laid on a likeness will grow a cause whether or not anybody put one there.

Mira went back down the yard the way she had come, with the smile still on the latch, and did not say she knew anything about the child under the fold.

\scenebreak

Nobody swept the \lore{Ash} early. A girl at the tubs started toward a broom when the last of the loaves went onto the rack, and Sella said, “Leave it. It'll be swept in the morning,” in the tone she used for \emph{don't set that near the door}, and the girl put the broom back. The woman at the tubs, who had said it before and would say it again tomorrow, said it now for the telling: the grey would be left where it lay through the Last Day, and burned in the morning, and called dealt with. She said it the way people say a debt has been settled when the slate has been turned down, as though a thing called dealt with had ever been a thing finished. Sella did not answer her. She went to the corner of the room where the grey ran along the skirting board and laid her thumb into it flat, the way she laid her hand on a child's belly to find whether it was warm.

It came to near two thumbs at Vail. She walked two doors down the \place{Clay Yard} and laid the same thumb into the ridge in the neighbour's room, where the grey ran from the bench by the door and turned the corner without a break; it came a finger's width thicker there. She walked back up the yard without saying which was which, and the ridge under the Vail skirting board stayed where it lay, turning its corner unbroken as it had turned it every morning for a year.

The boy lay in the basket at the edge of the boards and watched her lay the thumb down. He had learned to read the shape of a kept thing before he had learned to read anything else. A thing left in place was a thing still accounted for; a thing moved was a thing someone had decided about. Sella had not decided about the grey, and so the grey was still in the room, and the room was still answerable for it. He set the thumb beside the ridge on his private tray, and the ridge beside the charred slate lying in the kiln ash, and the slate beside the barrow of grey that went out through the north gate on market days, and did not join any of them. Four objects with one shape, and no order to put them in. Of \lore{the full note} there had been five hearings from four sources, and the fifth had come out of a dying man under a cold iron tube, and no one under this roof had added a sixth all day.

Five. From four sources. No sixth.

\scenebreak

The coat that had been \person{Garl Vail}'s kept its salt shape on its peg behind the door, and nobody took it down in all the baking. In \place{the Vail wheel-shed} the seat kept its hollow on the left, worn by twenty years of a man's leaning, and stayed hollow; the house had not filled it in three years and would not fill it now, on the Last Day, when every other place in the room was being put right for the turning of the year. The proverb went to the pot and to the fire and never to the basket: leave the crack showing; people trust a thing they can see has been handled. He had left the crack showing all year. The ridge in the corner was the crack he had left showing, and the house kept it in plain sight the way a surgeon keeps a wound dressed in a light room, so that nobody can say afterwards that they did not see it.

In the fold the two crusts lay side by side and were neither entered nor spoken of. One had been stiff since Stone 15 and one still smelled of barley smoke when the wind turned off \place{the Slip Brook}, though the wind had turned a hundred times since the smoke left it. A gift that cannot be answered cannot be refused; he had learned that in bread before he learned it in anything else, and had found it true of the fire, the basket, the wool, the warmed stone that was not his, and every other kindness in this house that had a payer who could be named and a recipient who could not thank. The women's cords and the house's loaves were the same kind of record: both kept without a name, both entered nowhere, both stubborn in the face of a slate. He could not have written either down in \place{Low Kiln}'s own hand if a slate and a stylus had been put into his palm, and he was the only person in the room who had seen both.

\person{Tarn} set the small knife back on its sill in the middle of the afternoon and told the yard where it was, unasked, in the tone of a man stating the weather. Six years of the seven were counted; the seventh he held back, as he held it back every day of the term, the one entry in this house that the house could not check. He scraped the grey from his knuckles with the flat of the knife, put the knife down again where it could be inspected, and went back to the broom. He had not said anything about the ridge being left, and he did not look at it now. The silence he kept was the one the whole house kept, and the two silences lay under the same roof without touching.

\person{Jorra Vail} came in from the shed with bone-coloured slip dried to old bone up both forearms and the bent little finger of his left hand going in ahead of him. Sella was telling him the fire was low; he moved the damper before she had finished the sentence, and said to the pot what he did not say to the room. “You'll do.” He did not say the back room was cold. He did not look toward the bed where the boy lay, and of the seat he meant to fill and the foundling he meant to give away he gave no sign at all. The last brooms stood dry against the wall, and the thatch above the door was dry enough to carry a spark if anybody had been careless with one, and nobody was.

At the bank she turned the coals with the back of a broken pot and did not touch \lore{the widow's mark} at her throat. A bank was not an opening, and the rule held on the Last Day as exactly as it held on any market's \lore{Account-day}. The coals went down under blind ash and the little-door burn on the back of her hand shut on nothing and gave nothing away. The kiln took its breath in fours and missed one, and the missing breath lay in the room beside the two crusts and the unswept ridge and the silence where a name should have stood, all of them kept, none of them entered, none of them joined.

The day ended with the brown-edged loaves cooling on the rack and the grey ridge running to the corner under the skirting board without a break, left where it lay to be burned in the morning and called gone. The boy lay twenty-nine days old and carried under the Vail roof, with the two crusts in the fold and the warmth entered nowhere. His hand, outside the wool, went up an inch toward where the fire had been and came down again.
